\chapter{Como a JVM funciona por dentro}
\label{cap:jvm-por-dentro}

\proximasaulas

O \texttt{javac} não gera código de máquina: gera bytecode, que a JVM
carrega, verifica, interpreta e, nos trechos quentes, compila. A memória
é gerenciada por um coletor de lixo que você escolhe, e o que uma thread
enxerga do que outra escreveu é definido pelo modelo de memória da
linguagem. Cada uma dessas peças muda uma decisão: quando medir, o que
deixar \texttt{static}, que GC usar, onde sincronizar.

As fontes primárias são a \citetitle{doc:jvms25} (que a partir daqui
chamamos de JVMS), a \citetitle{doc:jls25} (JLS) e a documentação da
HotSpot, a JVM do OpenJDK. O que é específico da HotSpot está marcado
como tal: a especificação permite outras implementações.

\section{Carregamento, linking e inicialização}
\label{sec:class-loading}

O cap.~5 da JVMS divide a vida de uma classe em três fases.

\begin{description}
\item[Carregamento (§5.3)] Um \termo{class loader}{Objeto que encontra os
  bytes de uma classe e a define na JVM.} encontra os bytes da classe.
  Há dois tipos: o \emph{bootstrap class loader}, da própria JVM, e os
  definidos pelo usuário, subclasses de \texttt{ClassLoader}. Um loader
  pode definir a classe ou delegar a outro.
\item[Linking (§5.4)] \emph{Verificação} confere que o bytecode é
  estruturalmente correto (falha: \texttt{VerifyError}).
  \emph{Preparação} cria os campos \texttt{static} com valores padrão,
  sem rodar código Java. \emph{Resolução} transforma referências
  simbólicas do \termo{constant pool}{Tabela de constantes do
  \texttt{.class}: nomes de classes, métodos, campos, strings e números.
  Ver a seção~\ref{sec:constant-pool}.} em referências concretas; a
  especificação permite fazer isso de forma preguiçosa, só quando a
  referência é usada.
\item[Inicialização (§5.5)] Roda o inicializador da classe
  (\texttt{<clinit>}: atribuições a campos \texttt{static} e blocos
  \texttt{static}). Acontece, por exemplo, na primeira execução de
  \texttt{new}, \texttt{getstatic}, \texttt{putstatic} ou
  \texttt{invokestatic} que referencia a classe.
\end{description}

\paragraph{O bug que isso explica.}
Se o bloco \texttt{static} de uma classe lança uma exceção (uma
configuração que não existe, um \texttt{Integer.parseInt} de variável de
ambiente vazia), o primeiro uso recebe \texttt{ExceptionInInitializerError}
com a causa original. A JVMS (§5.5, passos 11 e 12) manda marcar a classe
como \emph{erroneous}; a partir daí, todo uso recebe
\texttt{NoClassDefFoundError} (passo 5), que parece um problema de
classpath e não é. Quando vir \texttt{NoClassDefFoundError} num serviço
que funcionava, procure o primeiro \texttt{ExceptionInInitializerError}
no log.

\paragraph{A decisão que isso muda.}
Inicialização \texttt{static} pesada (abrir conexão, ler arquivo) roda no
primeiro uso, na thread que chegou primeiro, e transforma qualquer falha
nesse par de erros. Prefira inicialização explícita, num componente com
ciclo de vida (um bean), em que o erro aparece na subida.

\paragraph{Startup e o cache AOT.}
Carregar e ligar milhares de classes custa tempo a cada subida. O
\citetitle{doc:jep483} (JDK~24) guarda classes já carregadas e ligadas
num cache \termo{AOT}{\emph{Ahead-of-time}: compilação ou preparação
feita antes da execução, em oposição ao JIT, que compila durante a
execução.} gerado por uma execução de treino; o \citetitle{doc:jep514}
(JDK~25) reduz a criação do cache a um comando
(\texttt{java -XX:AOTCacheOutput=app.aot \ldots}), e o
\citetitle{doc:jep515} (JDK~25) guarda também perfis de métodos, para o
JIT começar mais cedo. O JEP~483 informa, sem dizer quem mediu nem em que
máquina, que o Spring PetClinic 3.2.0 subiu em 4,486~s no JDK~23 e em
2,604~s no JDK~24 com o cache (42\% mais rápido). Trate como número do
próprio projeto OpenJDK, não como garantia para o seu serviço.

\paragraph{Ponte com Quarkus.}
O guia \citetitle{doc:quarkus-class-loading} descreve uma hierarquia
própria em dev mode: um \emph{Base Runtime ClassLoader} persistente para
as dependências e um \emph{Runtime ClassLoader} para as classes da
aplicação, recriado quando a aplicação reinicia; é isso que permite o
live reload sem reiniciar a JVM. No pacote padrão de produção
(\texttt{fast-jar}), um loader que indexa as classes em tempo de build;
em executável nativo, classpath plano. A Parte~\ref{parte:quarkus} volta a isso.

\section{Interpretador e JIT}
\label{sec:jit}

A HotSpot começa interpretando o bytecode. Enquanto o programa roda, ela
mede onde o tempo é gasto e compila para código de máquina só o que é
quente; a documentação da Oracle descreve o compilador como
\emph{adaptativo} e diz que a maior parte do código, pouco usada, nem
chega a ser compilada (\citetitle{doc:jvm-overview}). Esse compilador é o
\termo{JIT}{\emph{Just-in-time}: compila durante a execução, com
informação sobre como o programa está de fato rodando.}.

\paragraph{Tiered compilation.}
Há dois compiladores, que o manual do comando \texttt{java}
(\citetitle{doc:java-cmd}) chama de C1 e C2. Com
\termo{tiered compilation}{Compilação em camadas: o método passa do
interpretador para o C1 (rápido de compilar, com coleta de perfil) e,
se continuar quente, para o C2 (mais lento de compilar, código mais
otimizado).}, introduzida no Java SE~7 e ligada por padrão, o
\emph{client compiler} (C1) gera versões que coletam perfil sobre si
mesmas, e o \emph{server compiler} (C2) usa esse perfil para compilar de
forma agressiva (\citetitle{doc:hotspot-perf}).

\paragraph{Interpretador, JIT ou AOT: quem decide, e quanto custa no startup?}
Como a JVM decide entre interpretador, JIT e AOT, e qual é o impacto real
do tempo de startup com compilação dinâmica e estática? Na HotSpot, quem
decide entre interpretador e JIT é a política de compilação, descrita no
comentário de cabeçalho de \citetitle{doc:hotspot-compilationpolicy}. Ela
tem cinco níveis: o 0 é o interpretador, o 1, o 2 e o 3 são o C1 com
graus diferentes de coleta de perfil, e o 4 é o C2, \enquote{with full
profile guided optimization}. A execução começa no nível~0, e a política
pode compilar o método no nível~3 ou no~2; terminada a coleta de perfil
no nível~3, o método passa ao nível~4. A passagem do nível~0 ao~3 compara
o número de invocações e de \emph{back-edges} (voltas de laço) do método
com limiares multiplicados por um coeficiente de escala, e o tamanho das
filas do C1 e do C2 ajusta esses limiares. Experimento (JDK~25.0.4.1,
imagem \texttt{eclipse-temurin:25-jdk}): \texttt{java
-XX:+PrintFlagsFinal -version} mostrou \texttt{Tier3InvocationThreshold}
igual a 200 e \texttt{Tier4InvocationThreshold} igual a 5000.

O AOT aparece de dois jeitos. Dentro da HotSpot, o
\citetitle{doc:jep483} (JDK~24) guarda num cache as classes carregadas e
ligadas (seção~\ref{sec:class-loading}), e o \citetitle{doc:jep515}
(JDK~25) guarda também perfis de métodos coletados num treino, para que o
JIT gere código nativo logo no início. Fora dela, o
\citetitle{doc:graalvm-native-image} compila o código Java, antes da
execução, para um executável nativo (seção~\ref{sec:jvm-vs-nativo}).

Os números publicados vêm dos próprios JEPs. O JEP~483 informa que um
programa pequeno roda em 0,031~s no JDK~23 e em 0,018~s no JDK~24 com o
cache, 42\% mais rápido; o número do Spring PetClinic está na
seção~\ref{sec:class-loading}. O JEP~515 informa que o seu programa de
exemplo roda em 90~ms com um cache sem perfis e em 73~ms depois de
coletar os perfis no cache, 19\% mais rápido. Para JVM contra executável
nativo, a seção~\ref{sec:jvm-vs-nativo} traz a tabela que o Quarkus
publica, com tempo até a primeira requisição, throughput de pico e
memória nos dois modos e com o cache AOT.

\paragraph{Otimizações que dependem do perfil.}
\begin{itemize}
\item \emph{Inlining}: o corpo do método chamado é copiado no chamador.
  É a porta para as outras otimizações, porque o compilador passa a
  enxergar os dois métodos juntos.
\item \termo{Escape analysis}{Análise que descobre se um objeto
  \enquote{escapa} do método ou da thread que o criou.}: se o objeto não
  escapa, o C2 pode eliminar a alocação (\emph{scalar replacement}) e os
  locks sobre ele (\emph{lock elision}). A documentação da Oracle
  ressalta que a HotSpot \emph{não} troca alocação no heap por alocação
  na pilha (\citetitle{doc:hotspot-perf}).
\item Especulação: o C2 compila só os caminhos que o perfil mostrou
  vivos e deixa nos outros uma \termo{uncommon trap}{Ponto no código
  compilado que devolve o controle ao runtime quando um caminho
  considerado impossível acontece.}. Se o caminho raro acontece, o
  código compilado é descartado
  (\termo{deoptimization}{Desotimização: trocar, em plena execução,
  frames de código otimizado por frames equivalentes de código não
  otimizado.}), o método volta
  ao interpretador ou a uma camada inferior, ganha novo perfil e é
  recompilado (\citetitle{artigo:shipilev-uncommon-traps}).
\item Chamadas virtuais: num ponto de chamada que só viu um tipo, o
  compilador pode \emph{devirtualizar}; com muitos tipos (megamórfico),
  a chamada vai por um despacho mais lento. Shipilëv mede a diferença e
  conclui que ela não justifica evitar chamadas virtuais
  (\citetitle{artigo:shipilev-megamorphic}).
\end{itemize}

\paragraph{Safepoints.}
Várias operações da JVM (algumas fases do GC, desotimização) precisam
parar as threads Java em pontos conhecidos, os
\termo{safepoints}{Pontos do código em que o estado da thread é
conhecido pela JVM; as threads verificam periodicamente se devem parar
ali.}. O código compilado tem verificações baratas (\emph{safepoint
polls}) que quase nunca disparam
(\citetitle{artigo:shipilev-safepoint}). O artigo explica por que os
polls existem: uma thread num laço ocupado pode nunca chamar a JVM por
conta própria, e por isso a HotSpot usa um esquema cooperativo, em que a
thread pergunta de tempos em tempos se deve parar num ponto de estado
conhecido. A opção \texttt{-Xlog:safepoint} registra os
safepoints no log da JVM.

\paragraph{As decisões que isso muda.}
\begin{itemize}
\item \emph{Warm-up é real.} As primeiras requisições rodam interpretadas
  ou em C1. Shipilëv avisa que o aquecimento precisa ter o
  mesmo perfil da produção (\citetitle{artigo:shipilev-uncommon-traps});
  senão, a recompilação aparece depois, no
  meio da medição. Teste de carga sem aquecimento mede o interpretador.
\item \emph{Microbenchmark à mão engana.} Um \texttt{System.nanoTime()}
  em volta de um laço mede compilação, eliminação de código morto e
  especulação. Use o \citetitle{doc:jmh}, cuja página pede que os
  benchmarks sejam revisados por pares.
\item \emph{Escreva código claro primeiro.} Métodos pequenos e objetos
  temporários que não escapam são exatamente o que o C2 otimiza melhor;
  o item 67 do \citetitle{livro:bloch2017} (\enquote{Optimize
  judiciously}) chega à mesma conclusão por outro caminho.
\end{itemize}

\paragraph{Vendo o JIT trabalhar.}
O manual do \texttt{java} documenta \texttt{-XX:+PrintCompilation}
(imprime uma linha a cada método compilado), \texttt{-Xint} (só
interpretador, para comparar), \texttt{-XX:-TieredCompilation} e, após
\texttt{-XX:+UnlockDiagnosticVMOptions}, \texttt{-XX:+PrintInlining} e
\texttt{-XX:+LogCompilation} (log em \texttt{hotspot\_pid<pid>.log} no
diretório atual; o manual do \texttt{java} ainda diz \texttt{hotspot.log},
e \texttt{-XX:LogFile} escolhe outro nome), que o
\citetitle{doc:jitwatch} visualiza). A seção~\ref{sec:ferramentas}
reúne essas ferramentas.

\section{Coleta de lixo}
\label{sec:gc}

Em Java você aloca e o \termo{garbage collector}{Coletor de lixo (GC):
componente da JVM que descobre quais objetos não são mais alcançáveis e
recupera a memória deles.} recupera. A escolha do coletor é uma troca
entre três coisas: \emph{throughput} (fração do tempo gasta no seu
código), latência (duração das pausas) e memória. O guia de GC da Oracle
para o JDK~25 (\citetitle{doc:gc-collectors}) descreve os coletores da
HotSpot assim:

\small
\begin{longtable}{@{}p{1.8cm} p{3.2cm} p{7.2cm}@{}}
\toprule
Coletor & Opção & Segundo a documentação \\
\midrule \endhead
Serial & \texttt{-XX:+UseSerialGC} & Uma thread faz todo o trabalho de
  GC; indicado para dados pequenos (até cerca de 100~MB) ou uma CPU \\
Parallel & \texttt{-XX:+UseParallelGC} & Como o Serial, mas com várias
  threads; para quando throughput máximo importa e pausas de um segundo
  ou mais são aceitáveis \\
G1 & \texttt{-XX:+UseG1GC} & Majoritariamente concorrente; busca cumprir
  uma meta de pausa com alta probabilidade mantendo bom throughput. É o
  padrão (\citetitle{doc:jep248}, JDK~9) \\
ZGC & \texttt{-XX:+UseZGC} & Baixa latência, pausas independentes do
  tamanho do heap; o guia fala em pausas máximas abaixo de um
  milissegundo, ao custo de algum throughput \\
\bottomrule
\end{longtable}
\normalsize

\paragraph{A evolução do ZGC.}
O ZGC ficou pronto para produção no \citetitle{doc:jep377} (JDK~15).
O \citetitle{doc:jep439} (JDK~21) acrescentou o modo geracional, o
\citetitle{doc:jep474} (JDK~23) o tornou padrão e o
\citetitle{doc:jep490} (JDK~24) removeu o modo não geracional. Em JDK~25,
\texttt{-XX:+UseZGC} é sempre o ZGC geracional; quem ainda passa
\texttt{ZGenerational} recebe aviso de opção obsoleta.

\paragraph{Shenandoah.}
O \citetitle{doc:jep379} (JDK~15) tornou o Shenandoah um coletor de
produção no OpenJDK; a página do projeto (\citetitle{doc:shenandoah}) o
descreve como um coletor de pausas ultrabaixas que faz mais trabalho
concorrente com o programa. Ele não aparece no guia de GC da Oracle
citado acima: confirme se a distribuição de JDK que você usa o inclui
antes de planejar com ele.

\paragraph{A decisão que isso muda.}
O guia da Oracle recomenda começar sem escolher nada: deixe a JVM
escolher, ajuste o tamanho do heap e só então troque de coletor, se o
requisito não for atendido. Na prática:
\begin{itemize}
\item API com requisito de latência (p99) apertado: compare G1 e ZGC
  medindo o p99 no teste de carga da Parte~I, não a média.
\item Job batch que só precisa terminar rápido: Parallel é candidato.
\item Container pequeno, com pouca memória e uma CPU: Serial é o que o
  guia indica para esse perfil.
\end{itemize}
Antes de trocar de GC, confira se o problema não é alocação demais: o
autoboxing da seção~\ref{sec:autoboxing} é um jeito comum de criar
objetos sem perceber. O JFR (seção~\ref{sec:ferramentas}) mostra pausas
e alocação por método.

\section{Modelo de memória}
\label{sec:modelo-memoria}

O capítulo~17 da JLS define o que uma thread pode ver das escritas de
outra. O conceito central é a relação
\termo{happens-before}{\enquote{Acontece antes}: se a escrita A
happens-before a leitura B, B enxerga o efeito de A. Sem essa relação
entre duas ações em threads diferentes, não há garantia de visibilidade
nem de ordem.} (§17.4.5). Exemplos da própria JLS: a escrita num campo
\texttt{volatile} happens-before toda leitura subsequente desse campo, e
um \texttt{unlock} de um monitor happens-before todo \texttt{lock}
subsequente do mesmo monitor. Campos \texttt{final} têm regras próprias
de publicação segura (§17.5).

Se um programa não tem \termo{data race}{Duas threads acessam a mesma
variável, pelo menos uma escreve, e não há happens-before entre os
acessos.}, ele se comporta como se fosse sequencial; com data races, uma
leitura pode ver valores surpreendentes. O texto de Shipilëv
(\citetitle{artigo:shipilev-jmm}) desenvolve essas consequências.

\paragraph{O bug que o bytecode mostra.}
Na listagem de \texttt{Sincronizacao} (seção~\ref{sec:sincronizacao}),
\texttt{contador++} ocupa os offsets 6 a 11: \texttt{getfield} lê o
campo, \texttt{iconst\_1} e \texttt{iadd} somam na pilha, e
\texttt{putfield} grava. São três passos, não um.
Entre o \texttt{getfield} e o \texttt{putfield}, outra thread pode ler o
mesmo valor antigo; as duas gravam o mesmo resultado e um incremento se
perde. \texttt{volatile} não resolve: ele garante visibilidade, mas o
ler-somar-gravar continua sendo três passos. As opções corretas são
\texttt{synchronized} (o que o exemplo faz), \texttt{AtomicInteger} ou
\texttt{LongAdder}, ou não compartilhar o estado. É o tema do item 78
do \citetitle{livro:bloch2017} (\enquote{Synchronize access to shared
mutable data}).

\paragraph{A decisão que isso muda.}
Records e campos \texttt{final} não são só estilo: objetos imutáveis
podem ser compartilhados entre threads sem lock (JLS §17.5 dá a garantia
para os campos \texttt{final}). Num serviço web, a pergunta \enquote{este
estado é compartilhado entre requisições?} vem antes de \enquote{qual
lock usar?}. O \citetitle{livro:goetz2006} é a referência para o resto.

\section{Onde aprofundar}

\subsection*{Livros}
\small
\begin{longtable}{@{}p{1.8cm} p{4.4cm} p{2.6cm} p{4cm}@{}}
\toprule
Tema & Livro & Autor(es), ano & Por que ler \\
\midrule \endhead
Desempenho & \citetitle{livro:oaks2020}, 2ª ed. & \citeauthor{livro:oaks2020},
  \citeyear{livro:oaks2020} & JIT, GC, memória, threads e metodologia de
  teste de desempenho, com foco em Java~8 e 11 \\
Desempenho & \citetitle{livro:evans2018} & \citeauthor{livro:evans2018},
  \citeyear{livro:evans2018} & Abordagem quantitativa: logs de GC, JIT,
  ferramentas e os perigos dos microbenchmarks \\
Concorrência & \citetitle{livro:goetz2006} & \citeauthor{livro:goetz2006},
  \citeyear{livro:goetz2006} & Visibilidade, publicação segura e
  utilitários de \texttt{java.util.concurrent} \\
\bottomrule
\end{longtable}
\normalsize

Os dois livros de desempenho são anteriores ao ZGC geracional e ao cache
AOT; use-os pelos conceitos e confira opções e padrões nos JEPs e na
documentação do JDK~25.

\subsection*{Artigos}
Da série \emph{JVM Anatomy Quarks}, de Aleksey Shipilëv:
\citetitle{artigo:shipilev-uncommon-traps},
\citetitle{artigo:shipilev-megamorphic} e
\citetitle{artigo:shipilev-safepoint}. Sobre o modelo de memória:
\citetitle{artigo:shipilev-jmm}.

\subsection*{Documentação}
\citetitle{doc:jvms25} (cap.~5), \citetitle{doc:jls25} (cap.~17),
\citetitle{doc:hotspot-perf}, \citetitle{doc:gc-collectors} e
\citetitle{doc:java-cmd}.

\section{Desafios práticos}

\desafio{Crie uma classe com um bloco \texttt{static} que lê uma
variável de ambiente com \texttt{Integer.parseInt}. Rode sem a variável
e chame a classe duas vezes, capturando a exceção. Quais exceções
aparecem em cada chamada? Relacione com os passos da JVMS §5.5.}

\desafio{Rode um programa seu com \texttt{-XX:+PrintCompilation} e
procure linhas com \enquote{made not entrant}. Que método foi
desotimizado? Que mudança no fluxo de entrada poderia ter causado isso?}

\desafio{Rode o mesmo teste de carga (Parte~\ref{parte:testes}) com
\texttt{-XX:+UseG1GC} e \texttt{-XX:+UseZGC}. Compare o p99 e o
throughput, e anote versão do JDK, tamanho do heap e máquina: sem isso,
o número não vale para mais ninguém.}

\desafio{Escreva uma classe com um \texttt{int} compartilhado e duas
threads que fazem \texttt{contador++} um milhão de vezes cada. Rode
algumas vezes. Depois troque por \texttt{volatile} e, por fim, por
\texttt{AtomicInteger}. Explique cada resultado usando o bytecode da
seção~\ref{sec:sincronizacao}.}

